\section{Wider material classes: where the template reaches and where it stops}\label{sec:wider}
The preceding sections follow one construction: a heat-flux kernel $G(s)$ multiplies the Fourier flux, and an effective modulus $m(s)$ carries the elastic side. This section places the other material classes of the literature relative to the template construction of \cref{sec:foundations}. What follows is a map of the literature rather than a derivation. The references listed here were matched to bibliographic records (authors, venue, volume, pages, DOI), but their equations were not re-derived in this review; statements below describe what each work is about and make no claim about its internal results. A class is called \emph{inside} the template if it enters through $G(s)$ or through a frequency-dependent modulus, and \emph{outside} if it changes the elastic operator, the list of fields, or the geometry.

\subsection{Couple-stress, strain-gradient and microstructured continua}
Higher-grade elasticity introduces an internal length into the elastic operator. The couple-stress theory was set out by Toupin \cite{Toupin1962,Toupin1964}, microstructure by Mindlin \cite{Mindlin1964} and the second gradient of strain by Mindlin \cite{Mindlin1965}; a multipolar formulation is due to Green and Rivlin \cite{GreenRivlin1964}. Plasticity-oriented gradient theories \cite{FleckHutchinson1993}, a couple-stress-based gradient theory \cite{Yang2002} and measurements interpreted with gradient elasticity \cite{Lam2003} show how the same internal length is fixed in different ways in different communities. Thermal extensions include thermo-microstretch solids \cite{Eringen1990ThermoMicrostretch}, thermoelastic damping of strain-gradient microplates \cite{Borjalilou2020}, Lord--Shulman thermoelasticity of microstructures \cite{Alihemmati2022}, nonlocal strain-gradient porous thermoelasticity \cite{Zeng2025}, Moore--Gibson--Thompson thermoelasticity under thermal shock \cite{HosseiniLi2022} and wave reflection at an interface with couple stress \cite{Li2018CoupleStress}. Wave problems in second-gradient media with damage \cite{Mawassy2023} and in pantographic metamaterials \cite{GodaGanghoffer2026} are recent examples from the mechanics side; gradient elasticity has also been used to remove dislocation singularities \cite{GutkinAifantis1999,LazarMaugin2005}.

\emph{Status relative to the template.} The elastic energy of these theories contains strain or rotation gradients, so the modulus acquires a wavenumber dependence \cite{Toupin1962,Mindlin1964,Mindlin1965}; this was read in the primary papers. That the kernel $G(s)$ itself is left unmodified is a structural statement, not a result of those papers. The isothermal papers contain no temperature, while the thermal versions of couple-stress and multipolar theory \cite{Toupin1962,GreenRivlin1964} carry their own heat equation and heat-flux law, which would have to be inserted into the kernel construction case by case. The length scale in this literature is normally fitted, and different enrichments can reproduce the same size effect; an internal length should therefore be quoted together with the experiment that fixed it, as in \cref{sec:size}.

\subsection{Phononic crystals and acoustic and thermal metamaterials}
Two mechanisms are usually distinguished. Bragg scattering in periodic composites was studied through the band structure of elastic and acoustic waves \cite{Sigalas1992,Kushwaha1993} (the latter reports a gap for transverse polarization only, complete throughout the Brillouin zone but not an absolute gap for all polarizations); local resonance gives gaps at wavelengths far larger than the lattice period \cite{Liu2000}, where the negative effective modulus near the resonance is the explanation offered in that paper, and surface-wave versions have been proposed \cite{Khelif2010}. Book-length treatments are \cite{Deymier2013,Laude2015}; the former defines the two gap mechanisms and prints a dynamic mass density for local resonance, the latter devotes a chapter to locally resonant crystals. Broad overviews are available \cite{Lu2009,Hussein2014,Christensen2015,Cummer2016,Ma2016,Vasileiadis2021,Cang2022}, as are topological edge-wave designs \cite{Wang2015Topo}. Studies that bring temperature or damping into such structures include thermal tuning of Lamb-wave band structure \cite{Yao2011}, enlargement of a locally resonant Lamb-wave gap \cite{Hu2014}, the influence of damping on transmission loss \cite{VanBelle2019} and a phononic-crystal strip for resonator quality factor \cite{Ha2022}. Heat-flux control is treated in the thermal diode \cite{Li2004Diode} and in the review of thermal metamaterials \cite{Li2021Thermal}.

\emph{Status relative to the template.} A locally resonant array is commonly described by a frequency-dependent effective modulus or mass (the original sonic-crystal paper \cite{Liu2000} uses an effective modulus; a negative-density example is given in \cite{Oh2016}), which is formally of the same kind as $m(s)$ in \cref{sec:modulus}; that is the point of contact. Coupling such an $m$ to the kernel $G(s)$ would give a thermoelastic damping estimate for a resonant array, but no reference located for this review treats that combination with the band structure included, so it is recorded as an open question in \cref{sec:open} and is not claimed as a result.

\subsection{Piezoelectric and piezoelectric-semiconductor waves}
Interaction of acoustic waves with carriers in piezoelectric semiconductors was described early \cite{Hutson1962} together with ultrasonic amplification \cite{White1962}, and surface-wave acoustoelectric interaction followed \cite{Palma1987}. Nanostructure applications span piezoelectric nanogenerators \cite{Wang2006Nano} and nanopiezotronics \cite{Wang2007Nano}; analyses of carrier distribution and electric potential in piezoelectric semiconductor structures \cite{Dai2018}, flexoelectric enhancement \cite{Zhao2022} and nonlinear thermo-deformation-polarization-carrier coupling \cite{Xiao2025} are recent. Flexoelectricity itself is reviewed in \cite{Cross2006,Zubko2013}. Photothermal and plasma-elastic-thermal wave studies in semiconductors \cite{Abbas2018,Khamis2020,Mahdy2020} form a thermoelastic branch of this field; a hyperbolic two-temperature treatment of a semiconductor with a cavity is in \cite{Hobiny2022}, and a nonlocal Kelvin--Voigt piezo-semiconductor with a memory-dependent heat flux, in which the flux law and the carrier equation are written together, is in \cite{Gupta2026HeatMap}, and temperature dependence of surface-wave characteristics has been measured \cite{Nishimura2005}. Quasicrystals, with phonon and phason fields, are a further extended-field class \cite{Shechtman1984,Levine1984,Chellappan2015,LiZhou2023}.

\emph{Status relative to the template.} Adding a carrier density, an electric potential or a phason field enlarges the list of unknowns and hence the dispersion polynomial, so these classes lie outside the single-kernel construction. A relaxation kernel for the heat flux has been inserted in the semiconductor models cited above \cite{Hobiny2022,Gupta2026HeatMap}; the early carrier-coupling papers \cite{Hutson1962,White1962} contain no heat equation, and the amplification criterion belongs to \cite{White1962}, not to \cite{Hutson1962}. No carrier equation is derived in this review.

\subsection{Graded, porous, laminated, magnetoelectric and nano materials; tissue}
Porous media go back to Biot \cite{Biot1941}; thermoelastic waves in saturated porous media \cite{Zorammuana2015}, micro-inertia effects \cite{SinghLianngenga2017} and nonlocal micropolar thermoelasticity \cite{Kalkal2020Refl}; the linear micropolar and microstretch thermal theory, with one temperature and one heat vector, is set out in \cite{Eringen1999}; these extend the picture. Functionally graded materials \cite{SureshMortensen1997} appear in thermoelastic damping of beams and sandwich structures \cite{Zhong2017,Peng2023}, thermal vibration and stability \cite{Fazzolari2018} and buckling of reinforced plates \cite{Meng2026}; temperature-dependent properties are surveyed in \cite{Noda1991}. Laminates are treated for thermoelastic damping in \cite{Vengallatore2005} and multifunctional composite mechanics in \cite{Gibson2010}. Magnetoelectric and multiferroic materials are reviewed in \cite{Eerenstein2006,Nan2008}. Carbon nanotubes and graphene \cite{Iijima1991,Lee2008} motivate the nano-scale beam and plate models of \cref{sec:size}, and laser-tissue interaction models \cite{Pearce2011,Li2026Bio} connect to the non-Fourier bioheat work of \cref{sec:applications}.

\emph{Status relative to the template.} Spatial grading, porosity and lamination mostly change coefficients or add interfaces, which the layered formulation of \cref{sec:layered} can accommodate; porous and magnetoelectric models add fields and are outside. The classification is by model structure and was not checked against each cited paper.

\subsection{Summary map}
\Cref{tab:wider} maps the wider material classes to their principal effects.
\begin{table}[ht]
\centering\small
\begin{tabularx}{\textwidth}{@{}>{\raggedright\arraybackslash}p{3.6cm}YY@{}}
\toprule
Class & What changes & Relative to the template \\
\midrule
Couple-stress, strain gradient & elastic operator, internal length & outside ($G$ unchanged, $m$ depends on $k$) \\
Dipolar, micropolar, microstretch & extra balance laws, larger elastic block & outside \\
Phononic and metamaterial & periodic structure; resonant effective mass or modulus & contact point through $m$; band structure outside \\
Piezoelectric semiconductors & carrier and potential fields & outside; heat-flux kernel can still be inserted \\
Quasicrystals & phonon and phason fields & outside \\
Graded, laminated & coefficients, interfaces & mostly inside, via the layered formulation \\
Porous, magnetoelectric & extra fields & outside \\
\bottomrule
\end{tabularx}
\caption{Material classes relative to the kernel construction. Entries describe model structure and are not results taken from the cited works.}\label{tab:wider}
\end{table}
